\chapter{2013年福建卷（文）}

\section{单选题}

\begin{problem}
复数 \(z = -1 - 2\i\)（\(\i\)为虚数单位）在复平面内对应的点位于（\quad）

\choices
    {第一象限}
    {第二象限}
    {第三象限}
    {第四象限}
\begin{answer}
C
\end{answer}

\begin{solution}
复数 \(z=-1-2\i\) 的实部为 \(-1\)，虚部为 \(-2\)，所以它在复平面内对应点的坐标为 \((-1,-2)\)，该点位于第三象限，故选 C.
\end{solution}
\end{problem}

\begin{problem}
设点 \(P(x, y)\)，则“\(x = 2\) 且 \(y = -1\)”是“点 \(P\) 在直线 \(l: x + y - 1 = 0\) 上”的（\quad）

\choices
    {充分而不必要条件}
    {必要而不充分条件}
    {充分必要条件}
    {既不充分也不必要条件}
\begin{answer}
A
\end{answer}

\begin{solution}
当 \(x=2\)，\(y=-1\) 时，\(x+y-1=2-1-1=0\)，所以该条件能推出点 \(P\) 在直线 \(l\) 上. 反过来，若点 \(P\) 在 \(l\) 上，只能得到 \(x+y=1\)，例如点 \((1,0)\) 也在 \(l\) 上，却不满足 \(x=2\) 且 \(y=-1\). 因此该条件是充分而不必要条件，故选 A.
\end{solution}
\end{problem}

\begin{problem}
若集合 \(A = \{1, 2, 3\}\)，\(B = \{1, 3, 4\}\)，则 \(A \cap B\) 的子集个数为（\quad）

\choices
    {\(2\)}
    {\(3\)}
    {\(4\)}
    {\(16\)}
\begin{answer}
C
\end{answer}

\begin{solution}
由 \(A=\{1,2,3\}\)，\(B=\{1,3,4\}\)，得 \(A\cap B=\{1,3\}\)，其中有 \(2\) 个元素. 含 \(2\) 个元素的集合共有 \(2^2=4\) 个子集，故选 C.
\end{solution}
\end{problem}

\begin{problem}
双曲线 \(x^{2} - y^{2} = 1\) 的顶点到其渐近线的距离等于（\quad）

\choices
    {\(\frac{1}{2}\)}
    {\( \frac{\sqrt{2}}{2} \)}
    {1}
    {\(\sqrt{2}\)}
\begin{answer}
B
\end{answer}

\begin{solution}
双曲线的顶点为 \((1,0)\) 和 \((-1,0)\)，其渐近线为 \(y=x\) 与 \(y=-x\). 以顶点 \((1,0)\) 为例，到直线 \(x-y=0\) 的距离为
\[
\frac{|1-0|}{\sqrt{1^2+(-1)^2}}=\frac{\sqrt{2}}{2}
\]
另一个顶点所得距离相同，故选 B.
\end{solution}
\end{problem}

\begin{problem}
函数 \( f(x) = \ln (x^2 + 1) \) 的图象大致是（\quad）

\choices
    {\choicebitmap[width=0.15\paperwidth]{img/2013/fujian_liberal/q5_optA.png}}
    {\choicebitmap[width=0.15\paperwidth]{img/2013/fujian_liberal/q5_optB.png}}
    {\choicebitmap[width=0.15\paperwidth]{img/2013/fujian_liberal/q5_optC.png}}
    {\choicebitmap[width=0.15\paperwidth]{img/2013/fujian_liberal/q5_optD.png}}
\begin{answer}
A
\end{answer}

\begin{solution}
函数的定义域为 \(\R\)，且 \(f(-x)=\ln((-x)^2+1)=f(x)\)，所以图象关于 \(y\) 轴对称. 又 \(x^2+1\geqslant1\)，从而 \(f(x)\geqslant0\)，且仅在 \(x=0\) 时取到最小值 \(0\). 此外，\(f'(x)=\dfrac{2x}{x^2+1}\)，所以 \(x>0\) 时函数递增，\(x<0\) 时函数递减. 因此图象是关于 \(y\) 轴对称、在原点处取最低点的曲线，对应选项 A.
\end{solution}
\end{problem}

\begin{problem}
若变量 \(x\)，\(y\) 满足约束条件 \(\begin{cases}
x + y \leqslant 2,\\
x \geqslant 1,\\
y \geqslant 0,
\end{cases}\) 则 \(z = 2x + y\) 的最大值和最小值分别为（\quad）

\choices
    {\(4\) 和 \(3\)}
    {\(4\) 和 \(2\)}
    {\(3\) 和 \(2\)}
    {\(2\) 和 \(0\)}
\begin{answer}
B
\end{answer}

\begin{solution}
约束区域的边界交点为 \((1,0)\)、\((2,0)\)、\((1,1)\). 在线性目标函数的三个顶点处计算：
\(z(1,0)=2\)，\(z(2,0)=4\)，\(z(1,1)=3\).
所以 \(z\) 的最大值为 \(4\)，最小值为 \(2\)，故选 B.
\end{solution}
\end{problem}

\begin{problem}
若 \(2^{x} + 2^{y} = 1\)，则 \(x + y\) 的取值范围是（\quad）

\choices
    {\([0,2]\)}
    {\([-2,0]\)}
    {\([-2, +\infty)\)}
    {\((- \infty, -2]\)}
\begin{answer}
D
\end{answer}

\begin{solution}
令 \(u=2^x\)，\(v=2^y\)，则 \(u>0\)，\(v>0\)，且 \(u+v=1\). 由均值不等式，\(0<uv\leqslant\frac14\)，其中上界在 \(u=v=\frac12\) 时取得；当 \(u\) 趋近于 \(0\) 或 \(1\) 时，\(uv\) 趋近于 \(0\)，但不取 \(0\). 因此
\[
x+y=\log_2(uv)\leqslant\log_2\frac14=-2
\]
且下方无界，故 \(x+y\in(-\infty,-2]\)，选 D.
\end{solution}
\end{problem}

\begin{problem}
阅读如图所示的程序框图，运行相应的程序，如果输入某个正整数 \(n\) 后，输出的 \(S \in (10, 20)\)，那么 \(n\) 的值为（\quad）

\bitmapfigure[max width=0.30\linewidth]{img/2013/fujian_liberal/q8_fig1.png}

\choices
    {\(3\)}
    {\(4\)}
    {\(5\)}
    {\(6\)}
\begin{answer}
B
\end{answer}

\begin{solution}
设循环执行 \(j\) 次. 初始 \(S=0\)，每次循环把 \(S\) 变为 \(1+2S\)，所以由递推可得 \(S=2^j-1\)；同时初始 \(k=1\)，每次循环后 \(k\) 加 \(1\)，故执行 \(j\) 次后 \(k=j+1\). 程序在 \(k>n\) 时停止，因此恰好执行 \(n\) 次，输出 \(S=2^n-1\). 由
\[
10<2^n-1<20
\]
得 \(11<2^n<21\)，所以 \(n=4\)，故选 B.
\end{solution}
\end{problem}

\begin{problem}
将函数 \( f(x) = \sin (2x + \theta)\)（\(-\frac{\pi}{2} < \theta < \frac{\pi}{2}\)）的图象向右平移 \(\varphi\)（\(\varphi > 0\)）个单位长度后得到函数 \( g(x) \) 的图象. 若 \(f(x)\)，\(g(x) \) 的图象都经过点 \( P\left(0, \frac{\sqrt{3}}{2}\right) \)，则 \(\varphi\) 的值可以是（\quad）

\choices
    {\(\frac{5\pi}{3}\)}
    {\( \frac{5\pi}{6} \)}
    {\( \frac{\pi}{2} \)}
    {\(\frac{\pi}{6}\)}
\begin{answer}
B
\end{answer}

\begin{solution}
函数图象向右平移 \(\varphi\) 个单位后，所得函数为 \(g(x)=f(x-\varphi)\). 因为 \(f(0)=\frac{\sqrt{3}}2\)，且 \(-\frac\pi2<\theta<\frac\pi2\)，所以 \(\sin\theta=\frac{\sqrt{3}}2\) 给出 \(\theta=\frac\pi3\). 又 \(g(0)=\frac{\sqrt{3}}2\)，故
\[
\sin\left(\frac\pi3-2\varphi\right)=\frac{\sqrt{3}}2
\]
因为 \(\sin t=\frac{\sqrt3}{2}\) 的解为 \(t=\frac\pi3+2k\pi\) 或 \(t=\frac{2\pi}3+2k\pi\)，取 \(\frac\pi3-2\varphi=\frac{2\pi}3-2\pi\)，得到 \(\varphi=\frac{5\pi}{6}\). 它满足 \(\varphi>0\)，故选 B.
\end{solution}
\end{problem}

\begin{problem}
在四边形 \(ABCD\) 中，\(\overrightarrow{AC} = (1, 2)\)，\(\overrightarrow{BD} = (-4, 2)\)，则该四边形的面积为（\quad）

\choices
    {\(\sqrt{5}\)}
    {\(2\sqrt{5}\)}
    {\(5\)}
    {\(10\)}
\begin{answer}
C
\end{answer}

\begin{solution}
四边形的面积等于两条对角线长度及其夹角正弦乘积的一半. 用向量坐标计算两条对角线的叉积：
\[
\left|\overrightarrow{AC}\times\overrightarrow{BD}\right|
=\left|1\cdot2-2\cdot(-4)\right|=10
\]
所以四边形 \(ABCD\) 的面积为 \(\frac12\times10=5\)，故选 C.
\end{solution}
\end{problem}

\begin{problem}
已知 \(x\) 与 \(y\) 之间的几组数据如下表：

\begin{center}
\begin{tabular}{c|cccccc}
\(x\) & \(1\) & \(2\) & \(3\) & \(4\) & \(5\) & \(6\) \\
\hline
\(y\) & \(0\) & \(2\) & \(1\) & \(3\) & \(3\) & \(4\)
\end{tabular}
\end{center}

假设根据上表数据所得线性回归直线方程为 \( \hat{y} = \hat{b}x + \hat{a} \)，若某同学根据上表中的前两组数据 \((1, 0)\) 和 \((2, 2)\) 求得的直线方程为 \( y = b'x + a' \)，则以下结论正确的是（\quad）

\choices
    {\(\hat{b} > b', \hat{a} > a'\)}
    {\(\hat{b} > b'\)，\(\hat{a} < a'\)}
    {\(\hat{b} < b'\)，\(\hat{a} > a'\)}
    {\(\hat{b} < b'\)，\(\hat{a} < a'\)}
\begin{answer}
C
\end{answer}

\begin{solution}
由表中数据得 \(\bar{x}=\frac{1+2+3+4+5+6}{6}=\frac72\)，\(\bar{y}=\frac{0+2+1+3+3+4}{6}=\frac{13}{6}\). 于是
\[
\sum_{i=1}^{6}(x_i-\bar{x})^2
=\left(-\frac52\right)^2+\left(-\frac32\right)^2+\left(-\frac12\right)^2+\left(\frac12\right)^2+\left(\frac32\right)^2+\left(\frac52\right)^2=\frac{35}{2}
\]
\[
\sum_{i=1}^{6}(x_i-\bar{x})(y_i-\bar{y})
=\frac{65+3+7+5+15+55}{12}=\frac{25}{2}
\]
因此回归直线的斜率与截距分别为
\(\hat b=\frac{25/2}{35/2}=\frac57\)，\(\hat a=\bar y-\hat b\bar x=\frac{13}{6}-\frac57\cdot\frac72=-\frac13\).
前两组数据确定的直线斜率、截距为 \(b'=\frac{2-0}{2-1}=2\)，\(a'=-2\). 于是 \(\hat b<b'\)，\(\hat a>a'\)，故选 C.
\end{solution}
\end{problem}

\begin{problem}
设函数 \( f(x) \) 的定义域为 \(\R\)，\(x_0 \) （\(x_0 \neq 0\)）是 \( f(x) \) 的极大值点，以下结论一定正确的是（\quad）

\choices
    {\(\forall x\in \R,f(x)\leqslant f(x_0)\)}
    {\(-x_0\) 是 \(f(-x)\) 的极小值点}
    {\(-x_0\) 是 \(-f(x)\) 的极小值点}
    {\(-x_0\) 是 \(-f(-x)\) 的极小值点}
\begin{answer}
D
\end{answer}

\begin{solution}
因为 \(x_0\) 是 \(f(x)\) 的极大值点，所以在 \(x_0\) 的某邻域内有 \(f(x)\leqslant f(x_0)\). 令 \(h(x)=-f(-x)\)，当 \(x\) 在 \(-x_0\) 的相应邻域内时，\(-x\) 在 \(x_0\) 的邻域内，从而
\[
h(x)=-f(-x)\geqslant-f(x_0)=h(-x_0)
\]
因此 \(-x_0\) 是 \(h(x)=-f(-x)\) 的极小值点，故选 D.
\end{solution}
\end{problem}

\section{填空题}

\begin{problem}
已知函数 \(f(x) = \begin{cases}
2x^{3}, & x < 0,\\
-\tan x, & 0 \leqslant x < \frac{\pi}{2},
\end{cases}\) 则 \(f\left(f\left(\frac{\pi}{4}\right)\right)=\fillinblank{}\).

\begin{answer}
\(-2\)
\end{answer}

\begin{solution}
因为 \(0\leqslant\frac\pi4<\frac\pi2\)，所以
\[
f\left(\frac\pi4\right)=-\tan\frac\pi4=-1
\]
又因为 \(-1<0\)，应使用分段函数的第一段，故
\[
f\left(f\left(\frac\pi4\right)\right)=f(-1)=2(-1)^3=-2
\]
\end{solution}
\end{problem}

\begin{problem}
利用计算机产生 \(0 \sim 1\) 之间的均匀随机数 \(a\)，则事件“\(3a - 1 < 0\)”发生的概率为\fillinblank{}.

\begin{answer}
\(\frac{1}{3}\)
\end{answer}

\begin{solution}
由 \(3a-1<0\) 得 \(0\leqslant a<\frac13\). 随机数 \(a\) 在区间 \([0,1]\) 上均匀分布，所以所求概率等于区间长度之比：
\[
P=\frac{\frac13-0}{1-0}=\frac13
\]
\end{solution}
\end{problem}

\begin{problem}
椭圆 \(\Gamma: \frac{x^2}{a^2} + \frac{y^2}{b^2} = 1\)（\(a > b > 0\)）的左右焦点分别为 \(F_1\)，\(F_2\)，焦距为 \(2c\). 若直线 \(y = \sqrt{3}(x + c)\) 与椭圆 \(\Gamma\) 的一个交点 \(M\) 满足 \(\angle MF_1F_2 = 2\angle MF_2F_1\)，则该椭圆的离心率等于\fillinblank{}.

\begin{answer}
\(\sqrt{3}-1\)
\end{answer}

\begin{solution}
取左、右焦点分别为 \(F_1=(-c,0)\)、\(F_2=(c,0)\). 直线 \(y=\sqrt3(x+c)\) 经过 \(F_1\)，其向右上方的射线与 \(x\) 轴正方向的夹角为 \(60^\circ\). 若点 \(M\) 在该直线的另一条射线上，则 \(\angle MF_1F_2=120^\circ\)，由题设得 \(\angle MF_2F_1=60^\circ\)，从而三角形在 \(M\) 处的内角为 \(0^\circ\)，与三角形非退化矛盾. 因此 \(M\) 在向右上方的射线上，从而
\[
\angle MF_1F_2=60^\circ
\]
由题设 \(\angle MF_1F_2=2\angle MF_2F_1\)，得 \(\angle MF_2F_1=30^\circ\)，于是三角形 \(MF_1F_2\) 在 \(M\) 处为直角三角形. 因为 \(F_1F_2=2c\)，所以由 \(30^\circ\)-\(60^\circ\)-\(90^\circ\) 三角形的边长关系，\(MF_1=c\)，\(MF_2=\sqrt3c\).

点 \(M\) 在椭圆上，故由椭圆定义
\[
2a=MF_1+MF_2=(1+\sqrt3)c
\]
因此离心率
\[
e=\frac ca=\frac{2}{1+\sqrt3}=\sqrt3-1
\]
\end{solution}
\end{problem}

\begin{problem}
设 \(S\)，\(T\) 是 \(\R\) 的两个非空子集. 如果存在一个从 \(S\) 到 \(T\) 的函数 \(y = f(x)\) 满足：
\begin{circlelist}
\item \(T = \{f(x)\mid x\in S\}\)；
\item 对任意 \(x_{1}\)，\(x_{2} \in S\)，当 \(x_{1} < x_{2}\) 时，恒有 \(f(x_{1}) < f(x_{2})\)，
\end{circlelist}
那么称这两个集合“保序同构”. 现给出以下 3 对集合：

\begin{circlelist}
\item \(A = \N\)，\(B = \N^{*};\)
\item \(A = \{x\mid -1\leqslant x\leqslant 3\}\)，\(B = \{x\mid -8\leqslant x\leqslant 10\}\)；
\item \( A = \{x \mid 0 < x < 1\} \)，\( B = \R \).
\end{circlelist}
其中，“保序同构”的集合对的序号是\fillinblank{}. （写出所有“保序同构”的集合对的序号）

\begin{answer}
\(\circled{1}\circled{2}\circled{3}\)
\end{answer}

\begin{solution}
对第\(\circled{1}\)对集合，按 \(\N=\{0,1,2,\ldots\}\)、\(\N^{*}=\{1,2,3,\ldots\}\) 取 \(f(x)=x+1\). 该函数严格递增，且值域为 \(\N^{*}\)，所以第\(\circled{1}\)对保序同构.

对第\(\circled{2}\)对集合，取
\[
f(x)=\frac92(x+1)-8=\frac92x-\frac72
\]
它把 \([-1,3]\) 的两个端点分别映为 \(-8,10\)，并且斜率为正，所以严格递增且值域为 \([-8,10]\)，第\(\circled{2}\)对保序同构.

对第\(\circled{3}\)对集合，取
\(f(x)=\tan\left(\pi x-\frac\pi2\right)\)，\(0<x<1\).
该函数在 \((0,1)\) 上严格递增，且当 \(x\to0^+\) 时 \(f(x)\to-\infty\)，当 \(x\to1^-\) 时 \(f(x)\to+\infty\)，值域为 \(\R\)，所以第\(\circled{3}\)对也保序同构. 因此应填 \(\circled{1}\circled{2}\circled{3}\).
\end{solution}
\end{problem}

\section{解答题}

\begin{problem}
已知等差数列 \(\{a_{n}\}\) 的公差 \(d = 1\)，前 \(n\) 项和为 \(S_{n}\).
\begin{enumerate}
\item 若 \(1\)，\(a_{1}\)，\(a_{3}\) 成等比数列，求 \(a_{1}\)；
\item 若 \(S_{5} > a_{1}a_{9}\)，求 \(a_{1}\) 的取值范围.
\end{enumerate}
\end{problem}

\begin{solution}
由等差数列的公差为 \(1\)，得
\(a_n=a_1+n-1\)，\(S_n=\frac n2\left(2a_1+n-1\right)\).
\begin{enumerate}
\item 由 \(1\)，\(a_1\)，\(a_3\) 成等比数列，得
\[
a_1^2=1\cdot a_3=a_1+2
\]
所以 \(a_1^2-a_1-2=0\)，即 \((a_1-2)(a_1+1)=0\)，从而 \(a_1=2\) 或 \(a_1=-1\).

\item 由 \(S_5>a_1a_9\)，且 \(S_5=5(a_1+2)=5a_1+10\)，\(a_9=a_1+8\)，得
\[
5a_1+10>a_1(a_1+8)
\]
即
\(a_1^2+3a_1-10<0\)，\((a_1+5)(a_1-2)<0\).
因此 \(a_1\) 的取值范围为 \((-5,2)\).
\end{enumerate}
\end{solution}

\begin{problem}
如图，在四棱锥 \(P-ABCD\) 中，\(PD \perp\) 平面 \(ABCD\)，\(AB \parallel DC\)，\(AB \perp AD\)，\(BC = 5\)，\(DC = 3\)，\(AD = 4\)，\(\angle PAD = 60^\circ\).
\begin{enumerate}
\item 当正视方向与向量 \( \overrightarrow{AD} \) 的方向相同时，画出四棱锥 \(P-ABCD\) 的正视图 （要求标出尺寸，并写出演算过程）；
\item 若 \(M\) 为 \(PA\) 的中点，求证： \(DM \parallel\) 平面 \(PBC\)；
\item 求三棱锥 \(D-PBC\) 的体积.
\end{enumerate}
\end{problem}

\begin{solution}
\begin{enumerate}
\item 在底面内取 \(D=(0,0,0)\)，使 \(AD\) 沿第一坐标轴，\(AB\) 沿第二坐标轴，并取 \(A=(4,0,0)\)，\(C=(0,3,0)\). 设 \(B=(4,t,0)\)，由 \(BC=5\) 得
\[
4^2+(t-3)^2=5^2
\]
由上式得 \(t=0\) 或 \(t=6\). 底面四边形非退化，\(t=0\) 时 \(B=A\)，故取 \(t=6\)，从而 \(AB=6\). 在直角三角形 \(PAD\) 中，
\[
PD=AD\tan\angle PAD=4\tan60^\circ=4\sqrt3
\]
故可取 \(P=(0,0,4\sqrt3)\).
所以沿 \(AD\) 方向正视时，\(A\)，\(D\) 投影重合，底边投影长度为 \(AB=6\)，竖直尺寸为 \(PD=4\sqrt3\). 正视图的外轮廓是顶点为 \(P\)、\(A(D)\)、\(B\) 的三角形，且 \(C\) 的投影位于底边中点，把底边分成 \(3\) 和 \(3\) 两段.

\item 由 \(M\) 为 \(PA\) 的中点，得 \(M=(2,0,2\sqrt3)\)，于是
\[
\overrightarrow{DM}=(2,0,2\sqrt3)
\]
又
\(\overrightarrow{PB}=(4,6,-4\sqrt3)\)，\(\overrightarrow{PC}=(0,3,-4\sqrt3)\)
并且
\[
\overrightarrow{DM}=\frac12\overrightarrow{PB}-\overrightarrow{PC}
\]
向量 \(\overrightarrow{PB}\)、\(\overrightarrow{PC}\) 是平面 \(PBC\) 内两条相交直线的方向向量，所以 \(\overrightarrow{DM}\) 是平面 \(PBC\) 的一个方向向量. 又
\[
\det\begin{pmatrix}
0&0&4\sqrt3\\
4&6&0\\
0&3&0
\end{pmatrix}=48\sqrt3\ne0
\]
故 \(D\notin\) 平面 \(PBC\)，从而 \(DM\parallel\) 平面 \(PBC\).

\item 以 \(D\) 为顶点，取 \(\overrightarrow{DP}\)、\(\overrightarrow{DB}\)、\(\overrightarrow{DC}\) 为棱，则
\[
V_{D-PBC}=\frac16\left|\det\begin{pmatrix}
0&0&4\sqrt3\\
4&6&0\\
0&3&0
\end{pmatrix}\right|
\]
行列式的绝对值为 \(48\sqrt3\)，所以
\[
V_{D-PBC}=\frac{48\sqrt3}{6}=8\sqrt3
\]
\end{enumerate}
\end{solution}

\begin{problem}
某工厂有 25 周岁以上（含 25 周岁）工人 300 名，25 周岁以下工人 200 名. 为研究工人的日平均生产量是否与年龄有关，现采用分层抽样的方法，从中抽取了 100 名工人，先统计了他们某月的日平均生产件数，然后按工人年龄在“25 周岁以上（含 25 周岁）”和“25 周岁以下”分为两组，再将两组工人的日平均生产件数分成 5 组：\([50, 60)\)，\([60, 70)\)，\([70, 80)\)，\([80, 90)\)，\([90, 100]\) 分别加以统计，得到如图所示的频率分布直方图.

\begin{enumerate}
\item 从样本中日平均生产件数不足 60 件的工人中随机抽取 2 人，求至少抽到一名“25 周岁以下组”工人的概率；
\item 规定日平均生产件数不少于 80 件者为“生产能手”，请你根据已知条件完成 \(2\times 2\) 列联表，并判断是否有 \(90\%\) 的把握认为“生产能手与工人所在的年龄组有关”.
\end{enumerate}
附：\(\chi^2 = \frac{n(ad - bc)^2}{(a + b)(c + d)(a + c)(b + d)}\)，\(P\left\{\chi^{2} \geqslant k\right\}\)

\begin{center}
\begin{tabular}{c|cccc}
\(P\{\chi^2 \geqslant k\}\) & \(0.100\) & \(0.050\) & \(0.010\) & \(0.001\) \\
\hline
\(k\) & \(2.706\) & \(3.841\) & \(6.635\) & \(10.828\)
\end{tabular}
\end{center}

（注：此处公式也可以写成 \(K^2 = \frac{n(ad - bc)^2}{(a + b)(c + d)(a + c)(b + d)}\)）

\begingroup
\leftskip=0pt
\rightskip=0pt
\bitmapfigure[max width=0.92\linewidth]{img/2013/fujian_liberal/q19_fig1.png}
\endgroup
\end{problem}

\begin{solution}
由分层抽样比例，样本中 \(25\) 周岁以上组有
\[
100\times\frac{300}{300+200}=60
\]
人，\(25\) 周岁以下组有 \(40\) 人. 每组直方图的组距均为 \(10\)，所以各组频率等于“频率组距”乘以 \(10\). 因此 \(25\) 周岁以上组在五个区间内的人数分别为 \(3\)，\(21\)，\(21\)，\(12\)，\(3\)，\(25\) 周岁以下组的人数分别为 \(2\)，\(10\)，\(13\)，\(13\)，\(2\).

\begin{enumerate}
\item 日平均生产件数不足 \(60\) 件的工人共有 \(3+2=5\) 人，其中 \(25\) 周岁以上组有 \(3\) 人. 随机抽取 \(2\) 人，至少抽到一名 \(25\) 周岁以下组工人的概率为
\[
1-\frac{\mathrm C_3^2}{\mathrm C_5^2}=1-\frac3{10}=\frac7{10}
\]

\item 日平均生产件数不少于 \(80\) 件的 \(25\) 周岁以上组工人有 \(12+3=15\) 人，\(25\) 周岁以下组工人有 \(13+2=15\) 人. 因此列联表为

\begin{center}
\begin{tabular}{c|ccc}
 & 生产能手 & 非生产能手 & 合计 \\
\hline
25 周岁以上组 & 15 & 45 & 60 \\
25 周岁以下组 & 15 & 25 & 40 \\
合计 & 30 & 70 & 100
\end{tabular}
\end{center}

由题给公式，取 \(a=15\)，\(b=45\)，\(c=15\)，\(d=25\)，\(n=100\)，得
\[
\chi^2=\frac{100(15\times25-45\times15)^2}{60\times40\times30\times70}=\frac{25}{14}\approx1.79
\]
因为 \(1.79<2.706\)，所以没有 \(90\%\) 的把握认为生产能手与工人所在的年龄组有关.
\end{enumerate}
\end{solution}

\begin{problem}
如图，抛物线 \(E: y^{2} = 4x\) 的焦点为 \(F\)，准线 \(l\) 与 \(x\) 轴的交点为 \(A\). 点 \(C\) 在抛物线 \(E\) 上. 以 \(C\) 为圆心，\(|CO|\) 为半径作圆，设圆 \(C\) 与准线 \(l\) 交于不同的两点 \(M\)，\(N\).
\begin{enumerate}
\item 若点 \(C\) 的纵坐标为 2，求 \( \left|MN\right| \)；
\item 若 \( \left|AF\right|^{2}=\left|AM\right|\cdot\left|AN\right| \)，求圆 \(C\) 的半径.
\end{enumerate}

\bitmapfigure[max width=0.72\linewidth]{img/2013/fujian_liberal/q20_fig1.png}
\end{problem}

\begin{solution}
抛物线 \(y^2=4x\) 的焦点为 \(F=(1,0)\)，准线为 \(l:x=-1\)，所以 \(A=(-1,0)\).

\begin{enumerate}
\item 当点 \(C\) 的纵坐标为 \(2\) 时，由抛物线方程得 \(C=(1,2)\)，圆 \(C\) 的半径为
\[
|CO|=\sqrt{1^2+2^2}=\sqrt5
\]
设圆与准线的交点为 \(M=(-1,y_1)\)、\(N=(-1,y_2)\)，则
\[
(-1-1)^2+(y-2)^2=5
\]
即 \((y-2)^2=1\)，所以 \(y_1=1\)，\(y_2=3\). 因此
\[
|MN|=3-1=2
\]

\item 设 \(C=(u,v)\)，则 \(v^2=4u\)，且圆的半径满足 \(r^2=|CO|^2=u^2+v^2=u^2+4u\). 圆与直线 \(x=-1\) 的交点纵坐标 \(y\) 满足
\[
(y-v)^2+(u+1)^2=u^2+4u
\]
从而
\[
(y-v)^2=2u-1
\]
因为交于不同的两点，令 \(s=\sqrt{2u-1}>0\)，则两交点的纵坐标为 \(v-s\)、\(v+s\). 又 \(|AF|=2\)，所以题设给出
\[
4=|AM|\cdot|AN|=|v-s|\cdot|v+s|=|v^2-s^2|
\]
由 \(v^2=4u\)、\(s^2=2u-1\)，有 \(v^2-s^2=2u+1>0\)，故 \(2u+1=4\)，即 \(u=\frac32\). 于是
\[
r=\sqrt{u^2+4u}=\sqrt{\frac94+6}=\frac{\sqrt{33}}2
\]
\end{enumerate}
\end{solution}

\begin{problem}
如图，在等腰直角 \(\triangle OPQ\) 中，\(\angle POQ = 90^{\circ}\)，\(OP = 2\sqrt{2}\)，点 \(M\) 在线段 \(PQ\) 上.
\begin{enumerate}
\item 若 \(OM = \sqrt{5}\)，求 \(PM\) 的长；
\item 若点 \(N\) 在线段 \(MQ\) 上，且 \(\angle MON = 30^{\circ}\)，问：当 \(\angle POM\) 取何值时，\(\triangle OMN\) 的面积最小？ 并求出面积的最小值.
\end{enumerate}

\FigureLayoutDeclare{img/2013/fujian_liberal/q21_fig1.png}{5.0cm}{7bp 71bp 101bp 3bp}
\bitmapfigure[max width=0.72\linewidth]{img/2013/fujian_liberal/q21_fig1.png}
\end{problem}

\begin{solution}
建立直角坐标系，取 \(O=(0,0)\)，\(P=(2\sqrt2,0)\)，\(Q=(0,2\sqrt2)\)，则直线 \(PQ\) 的方程为 \(x+y=2\sqrt2\).

\begin{enumerate}
\item 设 \(M=(x,y)\). 由 \(M\) 在 \(PQ\) 上及 \(OM=\sqrt5\)，得
\(x+y=2\sqrt2\)，\(x^2+y^2=5\).
于是
\[
2xy=(x+y)^2-(x^2+y^2)=8-5=3
\]
所以 \(x\)，\(y\) 是方程 \(t^2-2\sqrt2t+\frac32=0\) 的两个根，即 \(\{x,y\}=\left\{\frac{3\sqrt2}{2},\frac{\sqrt2}{2}\right\}\). 若 \(x=\frac{3\sqrt2}{2}\)，则
\[
PM=\sqrt{\left(\frac{3\sqrt2}{2}-2\sqrt2\right)^2+\left(\frac{\sqrt2}{2}\right)^2}=1;
\]
若 \(x=\frac{\sqrt2}{2}\)，则
\[
PM=\sqrt{\left(\frac{\sqrt2}{2}-2\sqrt2\right)^2+\left(\frac{3\sqrt2}{2}\right)^2}=3
\]
故 \(PM=1\) 或 \(3\).

\item 令 \(\theta=\angle POM\). 由于 \(M\)，\(N\) 均在 \(PQ\) 上且 \(N\) 在线段 \(MQ\) 上，\(0\leqslant\theta\leqslant60^\circ\)，并且 \(\angle PON=\theta+30^\circ\). 从 \(O\) 出发、与 \(OP\) 成角 \(\alpha\) 的射线与 \(PQ\) 的交点到 \(O\) 的距离为
\[
\frac{2\sqrt2}{\sin\alpha+\cos\alpha}
\]
因此
\(OM=\frac{2\sqrt2}{\sin\theta+\cos\theta}\)，\(ON=\frac{2\sqrt2}{\sin(\theta+30^\circ)+\cos(\theta+30^\circ)}\).
又 \(\angle MON=30^\circ\)，所以
\[
S_{\triangle OMN}=\frac12OM\cdot ON\sin30^\circ
=\frac{2}{(\sin\theta+\cos\theta)(\sin(\theta+30^\circ)+\cos(\theta+30^\circ))}
\]
记分母为 \(D\)，则
\[
D=2\sin(\theta+45^\circ)\sin(\theta+75^\circ)
=\cos30^\circ-\cos(2\theta+120^\circ)
\]
当 \(0\leqslant\theta\leqslant60^\circ\) 时，\(D\) 的最大值在 \(2\theta+120^\circ=180^\circ\)，即 \(\theta=30^\circ\) 时取得，且
\[
D_{\max}=1+\frac{\sqrt3}{2}
\]
所以三角形面积的最小值为
\[
\frac{2}{1+\frac{\sqrt3}{2}}=\frac{4}{2+\sqrt3}=8-4\sqrt3
\]
故当 \(\angle POM=30^\circ\) 时面积最小，最小值为 \(8-4\sqrt3\).
\end{enumerate}
\end{solution}

\begin{problem}
已知函数 \( f(x) = x - 1 + \frac{a}{\e^{x}} \)（\(a \in \R\)，\(\e\) 为自然对数的底数）.
\begin{enumerate}
\item 若曲线 \( y = f(x) \) 在点 \( (1, f(1)) \) 处的切线平行于 \( x \) 轴，求 \(a\) 的值；
\item 求函数 \( f(x) \) 的极值；
\item 当 \(a = 1\) 时，若直线 \(l: y = kx - 1\) 与曲线 \(y = f(x) \) 没有公共点，求 \(k\) 的最大值.
\end{enumerate}
\end{problem}

\begin{solution}
将函数写成 \(f(x)=x-1+a\e^{-x}\)，则
\[
f'(x)=1-a\e^{-x}
\]
\begin{enumerate}
\item 曲线在 \((1,f(1))\) 处的切线平行于 \(x\) 轴，等价于 \(f'(1)=0\)，所以
\(1-\frac a\e=0\)，\(a=\e\).

\item 当 \(a\leqslant0\) 时，\(f'(x)=1-a\e^{-x}>0\)，函数在 \(\R\) 上单调递增，没有极值.

当 \(a>0\) 时，\(f'(x)=0\) 等价于 \(\e^x=a\)，即 \(x=\ln a\). 当 \(x<\ln a\) 时 \(a\e^{-x}>1\)，所以 \(f'(x)<0\)；当 \(x>\ln a\) 时 \(a\e^{-x}<1\)，所以 \(f'(x)>0\). 因此函数在 \(x=\ln a\) 处取得极小值，且
\[
f(\ln a)=\ln a-1+\frac{a}{a}=\ln a
\]

\item 当 \(a=1\) 时，若直线 \(y=kx-1\) 与曲线有公共点，则存在实数 \(x\) 使
\(x-1+\e^{-x}=kx-1\)，\(\e^{-x}=(k-1)x\).
若 \(k>1\)，令 \(H(x)=\e^{-x}-(k-1)x\)，则 \(H(0)=1>0\)，且当 \(x\to+\infty\) 时 \(H(x)\to-\infty\)，由零点存在性定理，方程必有正根，直线必与曲线相交. 因此没有公共点时必有 \(k\leqslant1\). 当 \(k=1\) 时交点方程变为 \(\e^{-x}=0\)，无实数解，所以 \(k=1\) 确实满足没有公共点. 故 \(k\) 的最大值为 \(1\).
\end{enumerate}
\end{solution}
